% TOP_PROBLEM: {"id": 30006903, "title": "Rokhlin problem on multiple mixing", "category_id": 11, "category": {"id": 11, "name": "dynamical_systems", "display_name": "Dynamical Systems", "description": "Problems about long-term behavior of deterministic systems, Hamiltonian dynamics, and periodic orbits.", "slug": "dynamical-systems", "order_index": 11, "created_at": "2026-07-31T15:26:25.671Z"}, "status": "open"}
% FIRST_POSTED: 2026-09-25
% ENTRY_KIND: statement_only
% !TeX program = lualatex
% Definition and sources only; this file is not an LLM solution attempt.
\documentclass[11pt]{article}
\usepackage[margin=25mm]{geometry}
\usepackage{fontspec}
\setmainfont{Latin Modern Roman}
\usepackage{amsmath,amssymb,mathtools}
\usepackage{unicode-math}
\setmathfont{Latin Modern Math}
\usepackage{xurl,hyperref}
\hypersetup{colorlinks=true,urlcolor=blue}
\setcounter{secnumdepth}{0}
\setlength{\parindent}{0pt}
\setlength{\parskip}{5pt}
\setlength{\emergencystretch}{3em}
\begin{document}
\section*{\texorpdfstring{Rokhlin problem on multiple mixing}{Rokhlin problem on multiple mixing}}\label{30006903}
\subsection{Definitions and mathematical statement}
Let $T$ be an invertible measure-preserving transformation of a probability space $(X,\mu)$. Mixing means $\mu(A\cap T^{-n}B)\to\mu(A)\mu(B)$ as $|n|\to\infty$ for all measurable $A,B$. The question asks whether this implies, for every $k\ge3$ and measurable $A_1,\ldots,A_k$,
\[
 \mu\!\left(A_1\cap T^{-n_1}A_2\cap\cdots\cap
 T^{-(n_1+\cdots+n_{k-1})}A_k\right)
 \longrightarrow\prod_{j=1}^k\mu(A_j)
\]
as all positive gaps $n_i\to\infty$. Invertibility is the convention compatible with the catalogue's two-sided mixing definition; a one-sided version must be stated separately. The assertion is mixing of all orders, not simply decay of two-point correlations for selected observables, and does not impose a quantitative rate of convergence.
\par\smallskip\textit{Formulation and concept references:} \href{https://arxiv.org/abs/2411.07234}{[S1]}.

\subsection{Short English statement}
Does ordinary mixing of a probability-preserving transformation force statistical independence of every finite collection of observations separated by large time gaps?
\subsection{Sources}
\begin{itemize}
\item[S1]V. V. Ryzhikov, 'Multiple mixing, 75 years of Rokhlin's problem', arXiv:2411.07234 (2024).
\url{https://arxiv.org/abs/2411.07234}.
\textit{Formulation source.}
\end{itemize}
\end{document}
